% Section 1: Introductory Remarks
% Original pages: 347 (book page numbering)
% Translated from NT-Disk.pdf

\section{Introductory Remarks}
\label{sec:1}%
To analyze astrophysical aspects of black holes one must bring input from two
directions: from the elegant realm of general relativity (lectures of Hawking,
Carter, Bardeen, and Ruffini), and from the more mundane realm of astro-
physics and plasma physics (\S\,2 of these lectures). We shall here combine these
inputs to discuss the origin of stellar black holes (\S\,3); and to analyze observable
aspects of black holes in the interstellar medium (\S\,4), in binary star systems and
the nuclei of galaxies (\S\,5), and in cosmological contexts (\S\,6).

Readers with strong astrophysical backgrounds will wish to skip over our
brief treatment of the fundamentals of astrophysics and plasma physics (\S\,2),
and dig immediately into our discussion of black-hole astrophysics (\S\S\,3--8).
However, for the reader whose previous training has focussed primarily on
gravitation theory, the material of \S\,2 is an important prerequisite for under-
standing the rest of the notes.

The authors are deeply indebted to Mr.\ Alan Lightman for removing a large
number of errors from the manuscript.
